\subsection{ML-based network root cause diagnosis}

\ac{RCD} aims to identify the underlying reasons behind network anomalies, such as degraded performance or failures. The advent of machine learning (ML) has driven the development of numerous data-driven solutions for RCD, leveraging ML models to analyze network data and pinpoint the sources of issues. Below, we discuss notable research efforts in this domain, focusing on various supervised and unsupervised techniques tailored for network environments.

Dimopoulos et al. \cite{dimopoulos2015identifying} introduced a supervised ML framework to detect the root cause of QoE issues in video streaming. The metrics were collected from three key vantage points—mobile devices, routers, and content servers—to train supervised models. This work highlights the effectiveness of incorporating diverse data sources in diagnosing video streaming issues, enabling targeted interventions.
Kawasaki et al. \cite{kawasaki2020comparative} focused on root cause analysis in \ac{NFV} environments. The study applied supervised ML algorithms such as MLP, SVM, and Random Forests for fault classification. By comparing multiple ML models, the authors demonstrate the feasibility of using supervised approaches for diagnosing issues in complex, virtualized networks.
Salik et al. \cite{salik2023non} detected non-WiFi interference and estimated its impact on WiFi throughput for both 2.4GHz and 5GHz bands. The authors developed supervised models, including Random Forest, CatBoost, and MLP, using network metrics collected at the WiFi edge. This approach effectively identified and mitigated interference from non-WiFi sources, improving network reliability in congested environments.
Syrigos et al. \cite{syrigos2019employment} developed ML-based tools to troubleshoot common WiFi impairments such as contention, interference, hidden terminals, and low signal strength. Supervised ML models were trained on curated datasets to classify network impairments, enabling actionable diagnostics. The framework provided a systematic way to enhance WiFi network performance by pinpointing specific issues.
Dötterl et al. \cite{dotterl2024classification} classified home network problems using advanced deep learning techniques. A transformer-based architecture is employed to analyze outputs from network monitoring tools like ping and dig. It demonstrated that transformer models can effectively classify complex network problems, offering precise diagnostics for home networks.

Fida et al. \cite{fida2023bottleneck} proposed a two-stage framework for bottleneck identification in cloudified 5G networks: in Stage 1, an unsupervised \ac{VAE} detects anomalies in network telemetry data while in Stage 2, an MLP classifier identifies the specific root cause of the bottlenecks. The distributed telemetry framework enabled real-time diagnosis in 5G networks, demonstrating scalability and adaptability in modern mobile environments.

Our work builds on recent advancements in \ac{RCD} by addressing the limitations of existing solutions that rely exclusively on supervised methods. We introduce a novel approach specifically tailored for low-latency applications operating in Wi-Fi network environments. By leveraging contrastive learning, our solution enhances diagnostic accuracy, particularly in scenarios with limited labeled data. While our approach shares similarities with the method proposed in \cite{fida2023bottleneck}, we differentiate ourselves by employing contrastive learning for anomaly detection, which results in improved performance.

% Zhao et al. \cite{zhao2024fault} introduced a graph-based convolutional network (GCN) method for 5G networks, leveraging user-side KPIs and synthetic data generated by GANs. Their approach also incorporated expert-based correlation graphs to enhance model training and interpretability

% Zhang et al. \cite{zhang2018wifi} WiFi and Multiple Interfaces: Adequate for Virtual Reality? 
% Burbano et al. \cite{burbano2024impact} Impact of Network Factors on QoE-sensitive Virtual Reality Environments: An Experimental Analysis Perspective using VR-based Mindfulness

\subsection{Time series Classification}
\ac{TSC} has been a prominent research topic for decades due to its applicability in numerous real-world domains, including cybersecurity and network management. A variety of solutions have been proposed in the literature, with methods ranging from classical approaches to modern deep learning techniques.

According to Bagnall et al. \cite{bagnall2017great}, classical TSC methods can be categorized into several groups including distance-based methods amoung those one of the most prominent examples is the 1-NN-DTW algorithm, which performs one-nearest-neighbor classification using dynamic time warping (DTW) as the distance metric; interval-based methods like Time Series Forest (TSF) \cite{deng2013time} which leverage random forests and extract summary statistics from intervals of time series data to improve classification accuracy; shapelet-based methods like binary shapelet transformation methods \cite{bostrom2015binary} which use shapelets that are phase-independent subseries that capture discriminative patterns within time series data for TSC or ensemble-based methods like Collective of Transformation-based Ensembles (COTE) \cite{bagnall2015time} that combines classifiers trained on different data representations. This approach highlights the importance of transforming data into subspaces where features are more easily discriminable.

In recent years, deep learning techniques have gained significant traction in TSC due to their ability to learn complex and hierarchical features directly from raw data. Various architectures have been proposed, including MLPs, CNNs, TCNs or RNNs \cite{zheng2016exploiting, wang2017time, zhao2017convolutional}. These techniques were supervised learning techniques and were successful in time series classification over diverse domains \cite{ismail2019deep}. However, the rise of SSL have come with a jump in performance of TSC models. In fact, SSL has gained significant interest for its ability to learn meaningful representations from unlabeled data through carefully designed pretext tasks on computer vision or natural language processing. These tasks aim to uncover the intrinsic structure of the data, enabling the learned representations to be effectively utilized in downstream tasks, including TSC. Among the various SSL approaches, CL has emerged as particularly effective. It works by encouraging similar samples (positive pairs) to cluster in the feature space while pushing apart dissimilar samples (negative pairs). For time series data, numerous CL frameworks have been proposed, demonstrating strong performance on TSC tasks. T-Loss \cite{franceschi2019unsupervised} introduces a triplet loss that constructs positive pairs using subseries extracted from the same time series, while negative pairs are formed using subseries from other time series. Coupled with a SVM, the learned representations perform exceptionally well on multivariate TSC tasks, demonstrating robust feature extraction. TNC (Temporal Neighborhood Coding) \cite{tonekaboni2021unsupervised} constructs pairs based on the concept of temporal neighborhoods, leveraging statistical tests to identify segments of a time series that share similar underlying properties. It achieves excellent results on simulated and real-world datasets, such as ECG data, by capturing temporal dependencies. TS-TCC (Time-Series Temporal and Contextual Contrastive Learning) \cite{Eldele2021TimeSeriesRL} employs cross-view prediction using data augmentations and integrates temporal and contextual contrastive learning modules. This design ensures the model learns robust, discriminative representations and demonstrates superior accuracy on TSC tasks across diverse domains by effectively handling temporal dynamics. CLOCS \cite{kiyasseh2021clocs} specifically designed for ECG signals, applies data transformations to learn patient-specific representations that are invariant to spatio-temporal variations. Focused on healthcare, this method excels in classifying ECG time series by capturing unique and invariant features. TF-C (Time-Frequency Contrastive Framework) \cite{zhang2022self} encourages time series with time and frequency similarity to cluster in the time-frequency space thanks to the introduction of novel frequency-based augmentation techniques to generate augmented views of the data. It achieves competitive results in TSC across a wide range of domains, showcasing the effectiveness of incorporating frequency-domain information. TS2Vec \cite{Yue2021TS2VecTU} proposes a hierarchical framework to learn both temporal-wise and instance-wise similarities at multiple semantic levels. The contextual consistency-based pair selection enhances representation learning for diverse time series patterns. It demonstrates versatility and efficiency in multivariate TSC and other related tasks, further establishing the value of contextual representation learning.

We tackle RCD in Cloud VR as a time series classification problem. Unlike aforementioned approaches to time series classification, our method employs a two-stage framework: first, it identifies normal scenario time series, and then it classifies faulty scenarios, thereby improving the overall accuracy and effectiveness of the classification process.


 

% \subsection{Anomaly detection techniques}

% Anomaly detection (AD) has been extensively studied, with methods broadly categorized into machine learning and deep learning based approaches. Traditional statistical methods assume that anomalies deviate from high-probability regions of a probabilistic model. Parametric methods like ARIMA \cite{yaacob2010arima} and Gaussian-based techniques \cite{luo2017gas} rely on specific distributions, while non-parametric methods, such as \ac{KDE} \cite{cao2016one}, estimate density functions without such assumptions. Distance-based methods, including k-Nearest Neighbors (kNN) \cite{chaovalitwongse2007time} and \ac{LOF} \cite{breunig2000lof}, identify anomalies by measuring distances or densities in feature space. Additionally, spectral analysis methods like Principal Component Analysis (PCA) \cite{paffenroth2018robust} and isolation-based techniques like Isolation Forest (IF) \cite{liu2008isolation} are commonly used for AD.  

% Deep learning methods have significantly advanced AD, particularly in complex time series and high-dimensional data. Reconstruction-based techniques, such as autoencoders (AEs), detect anomalies by measuring reconstruction errors. Advanced models like VAEs \cite{xu2018unsupervised} and DAGMM \cite{zong2018deep} combine probabilistic models and dimensionality reduction for improved performance. Temporal dependencies are captured using RNN-based methods like LSTM-VAE \cite{park2018multimodal} and Omni-Anomaly \cite{su2019robust}, while Transformer-based models such as Anomaly Transformer \cite{xu2022anomaly} and TranAD \cite{tuli2022tranad} leverage attention mechanisms for long-sequence AD. Generative models, including GAN-based frameworks like MAD-GAN \cite{li2019mad} and BeatGAN \cite{zhou2019beatgan}, focus on reconstructing time series data while adversarially training models to enhance robustness. Forecasting-based methods, such as LSTM-PRED \cite{goh2017anomaly} and Telemanom \cite{hundman2018detecting}, use deviations between predicted and observed values for anomaly identification.  

% Recent advancements in contrastive learning have further improved AD performance by enhancing representation learning. Frameworks like COCA \cite{wang2023deep} and ContrastAD \cite{li2023contrastive} leverage contrastive losses to train models on positive and negative pairs, with techniques like data augmentation and synthetic anomaly generation enriching the training process. CARLA \cite{DARBAN2025110874} combines self-supervised learning with contrastive loss to distinguish normal and anomalous samples effectively, while DCdetector \cite{yang2023dcdetector} employs multi-scale attention mechanisms to capture spatial and temporal dependencies. These approaches highlight the evolution of AD from traditional statistical methods to advanced DL techniques, emphasizing the growing importance of representation learning and model adaptability in diverse application scenarios.  

% Our work leverages anomaly detection techniques as a foundational step in our two-stage framework for RCD, aiming to improve classification accuracy. Similar to prior studies, we incorporate contrastive learning to enhance the representation of both anomalous and normal time series. This approach not only strengthens the detection of anomalies but also significantly improves the performance of the subsequent classification task, setting our method apart in effectively addressing RCD challenges.
